\chapter{Scalability Blueprint}
\label{ch:scalability}

This chapter details the architectural design and roadmap for scaling RGUKT Connect to support high user traffic and concurrent messaging.

\section{Horizontal Scaling of Backend Pods}
The Spring Boot backend is packaged as a stateless Docker container. This allows the application to be scaled horizontally by running multiple replicas in the Kubernetes cluster. Traffic is distributed across these replicas using a Kubernetes Service.

However, scaling the stateless API layer introduces challenges for features that manage state, such as WebSockets and in-memory caches.

\section{Redis Message Broker for WebSocket Sync}
Currently, the WebSocket STOMP broker is run in-memory within each backend process. If a client connects to Pod A, and another connects to Pod B, they cannot exchange messages directly because their sessions are isolated on different servers.

To support multiple backend pods, the application must use a shared message broker (such as RabbitMQ or a Redis Pub/Sub backplane) instead of the in-memory broker.

\begin{figure}[H]
    \centering
    \begin{tikzpicture}[node distance=1.5cm, auto]
        \node [block, fill=LightGrey] (clientA) {Client A};
        \node [block, right=4cm of clientA, fill=LightGrey] (clientB) {Client B};
        
        \node [block, below=1cm of clientA] (podA) {Backend Pod A};
        \node [block, below=1cm of clientB] (podB) {Backend Pod B};
        
        \node [block, below=1.5cm of $(podA.south)!0.5!(podB.south)$, fill=RguktMaroon, text=white, minimum width=8cm] (redis) {Redis Pub/Sub / STOMP Broker Backplane};

        \draw [arrow] (clientA) -- (podA);
        \draw [arrow] (clientB) -- (podB);
        
        \draw [arrow, <->] (podA) -- (redis -| podA);
        \draw [arrow, <->] (podB) -- (redis -| podB);
    \end{tikzpicture}
    \caption{Distributed WebSocket Architecture using a Redis Pub/Sub Backplane}
    \label{fig:redis_pubsub}
\end{figure}

As shown in Figure~\ref{fig:redis_pubsub}, when Pod A receives a message, it publishes it to the Redis channel. Pod B, which is subscribed to the Redis channel, receives the message and forwards it to Client B, enabling cross-pod messaging.

\section{Database Scaling: Read-Write Separation}
Because user activity in social platforms is read-heavy (e.g., browsing feeds, searching directories, checking jobs), the database can become a bottleneck. The following database architecture is recommended for scaling:
\begin{itemize}
    \item \textbf{Write Master Database:} A single database instance that processes all write operations (e.g., new posts, connections, chat logs).
    \item \textbf{Read Replicas:} Multiple read-only databases synced asynchronously with the master. The application routes read operations (such as directory queries) to these replicas, distributing the database load.
\end{itemize}

\section{Content Delivery Network (CDN) Integration}
To reduce load on the origin server, static assets (such as profile images and post attachments uploaded to Amazon S3) should be served through a CDN (such as Amazon CloudFront or Cloudflare). The CDN caches these assets at edge locations close to users, reducing latency and bandwidth costs.
